% GENERATED FILE. Edit canonical lessons, then run npm run generate:books.
% canonical-source-hash: fnv1a64:41ef625b79be46d7
% canonical-lessons: JA-C127-ooi, JA-C127-sukunai, JA-C127-sawayaka, JA-C127-odayaka, JA-C127-sunao

\chapter{Many, Few, Refreshing, Calm, Honest}
\label{ch:ja-many-few-refreshing-calm-honest}
\hlchaptermodality{127}

\begin{chapteropening}
\textbf{By the end of this chapter:} \emph{I can say many, few, refreshing, calm and honest.}

Complete the last lesson of chapter 127: I can say many, few, refreshing, calm and honest.
\end{chapteropening}

\section[\texorpdfstring{\ja{おおい} (ooi)}{おおい (ooi)}]{\ja{おおい} (ooi) — many, a lot}

\label{lesson:JA-C127-ooi}



\begin{quote}
Before the new one: say the Japanese for ordinary, then the Japanese for tasty; skilful.
\end{quote}

\subsection*{You'll want to know: \ja{おおい}}
\textbf{\ja{おおい}} — \emph{ooi} — \textquotedblleft{}many, a lot\textquotedblright{}.

\begin{grammarlens}[title={What you've built}]
One more word for this chapter.
\end{grammarlens}

\subsection*{Your turn}
\begin{itemize}
\raggedright
  \item \emph{Say it:} \emph{ooi}
  \item \emph{Say it:} \emph{ooi}, once more
  \item \emph{Say it:} say \emph{ooi}
  \item \emph{Recall:} say the Japanese for ordinary, then the Japanese for tasty; skilful, then say \emph{ooi} again
\end{itemize}

\begin{tcolorbox}[breakable,colback=teal!4,colframe=teal!35!black,title={Before you move on}]
What does \ja{おおい} mean? (\textquotedblleft{}Many, a lot\textquotedblright{}.) Say it once more.
\end{tcolorbox}
\section[\texorpdfstring{\ja{すくない} (sukunai)}{すくない (sukunai)}]{\ja{すくない} (sukunai) — few, little}

\label{lesson:JA-C127-sukunai}



\begin{quote}
Before the new one: say the Japanese for tasty; skilful, then the Japanese for many.
\end{quote}

\subsection*{You'll want to know: \ja{すくない}}
\textbf{\ja{すくない}} — \emph{sukunai} — \textquotedblleft{}few, little\textquotedblright{}.

\begin{grammarlens}[title={What you've built}]
One more word for this chapter.
\end{grammarlens}

\subsection*{Your turn}
\begin{itemize}
\raggedright
  \item \emph{Say it:} \emph{sukunai}
  \item \emph{Say it:} \emph{sukunai}, once more
  \item \emph{Say it:} say \emph{sukunai}
  \item \emph{Recall:} say the Japanese for tasty; skilful, then the Japanese for many, then say \emph{sukunai} again
\end{itemize}

\begin{tcolorbox}[breakable,colback=teal!4,colframe=teal!35!black,title={Before you move on}]
What does \ja{すくない} mean? (\textquotedblleft{}Few, little\textquotedblright{}.) Say it once more.
\end{tcolorbox}
\section[\texorpdfstring{\ja{さわやか} (sawayaka)}{さわやか (sawayaka)}]{\ja{さわやか} (sawayaka) — refreshing}

\label{lesson:JA-C127-sawayaka}



\begin{quote}
Before the new one: say the Japanese for many, then the Japanese for few.
\end{quote}

\subsection*{You'll want to know: \ja{さわやか}}
\textbf{\ja{さわやか}} — \emph{sawayaka} — \textquotedblleft{}refreshing\textquotedblright{}.

\begin{grammarlens}[title={What you've built}]
One more word for this chapter.
\end{grammarlens}

\subsection*{Your turn}
\begin{itemize}
\raggedright
  \item \emph{Say it:} \emph{sawayaka}
  \item \emph{Say it:} \emph{sawayaka}, once more
  \item \emph{Say it:} say \emph{sawayaka}
  \item \emph{Recall:} say the Japanese for many, then the Japanese for few, then say \emph{sawayaka} again
\end{itemize}

\begin{tcolorbox}[breakable,colback=teal!4,colframe=teal!35!black,title={Before you move on}]
What does \ja{さわやか} mean? (\textquotedblleft{}Refreshing\textquotedblright{}.) Say it once more.
\end{tcolorbox}
\section[\texorpdfstring{\ja{おだやか} (odayaka)}{おだやか (odayaka)}]{\ja{おだやか} (odayaka) — calm}

\label{lesson:JA-C127-odayaka}



\begin{quote}
Before the new one: say the Japanese for few, then the Japanese for refreshing.
\end{quote}

\subsection*{You'll want to know: \ja{おだやか}}
\textbf{\ja{おだやか}} — \emph{odayaka} — \textquotedblleft{}calm\textquotedblright{}.

\begin{grammarlens}[title={What you've built}]
One more word for this chapter.
\end{grammarlens}

\subsection*{Your turn}
\begin{itemize}
\raggedright
  \item \emph{Say it:} \emph{odayaka}
  \item \emph{Say it:} \emph{odayaka}, once more
  \item \emph{Say it:} say \emph{odayaka}
  \item \emph{Recall:} say the Japanese for few, then the Japanese for refreshing, then say \emph{odayaka} again
\end{itemize}

\begin{tcolorbox}[breakable,colback=teal!4,colframe=teal!35!black,title={Before you move on}]
What does \ja{おだやか} mean? (\textquotedblleft{}Calm\textquotedblright{}.) Say it once more.
\end{tcolorbox}
\section[\texorpdfstring{\ja{すなお} (sunao)}{すなお (sunao)}]{\ja{すなお} (sunao) — honest, straightforward}

\label{lesson:JA-C127-sunao}



\begin{quote}
Before the new one: say the Japanese for refreshing, then the Japanese for calm.
\end{quote}

\subsection*{You'll want to know: \ja{すなお}}
\textbf{\ja{すなお}} — \emph{sunao} — \textquotedblleft{}honest, straightforward\textquotedblright{}.

\begin{grammarlens}[title={What you've built}]
One more word for this chapter.
\end{grammarlens}

\subsection*{Your turn}
\begin{itemize}
\raggedright
  \item \emph{Say it:} \emph{sunao}
  \item \emph{Say it:} \emph{sunao}, once more
  \item \emph{Say it:} say \emph{sunao}
  \item \emph{Recall:} say the Japanese for refreshing, then the Japanese for calm, then say \emph{sunao} again
\end{itemize}

\begin{tcolorbox}[breakable,colback=teal!4,colframe=teal!35!black,title={Before you move on}]
What does \ja{すなお} mean? (\textquotedblleft{}Honest, straightforward\textquotedblright{}.) Say it once more.
\end{tcolorbox}
